%! Unit 1: Vectors and Matrices — Summative Assessment — Practice Test --- Study Copy (KEY)
%
% THE WORK RULE: every \begin{work} block below is authored byte-identically in
% ../../tests/practice_test/main.tex. There it ships as a \vphantom of the same
% height; here linalg-key makes it visible in keyred. Edit a work block in BOTH
% files or in neither — that is what keeps this key the same length as its blank.
%
% Teacher prose (Part B rationale, Part D scoring) is NOT here: it lives on page 2
% of ../../unit_cover_key/main.tex, which reaches the key packet only.
\documentclass[10pt]{article}
\usepackage{linalg-article}
\usepackage{linalg-key}   % pulls in -boxes; do NOT also load -boxes
\newcommand{\vv}{\mathbf{v}}
\newcommand{\ww}{\mathbf{w}}
\newcommand{\uu}{\mathbf{u}}
\newcommand{\bb}{\mathbf{b}}
\newcommand{\xx}{\mathbf{x}}

% Part-divider strip
\newcommand{\parthead}[1]{%
  \vspace{6pt}\noindent
  \begin{headlinebox}{burgundy}{\color{white}\bfseries #1}\end{headlinebox}%
  \vspace{4pt}}

\begin{document}

\pageheader{Unit 1: Vectors and Matrices}{Practice Test --- Study Copy --- Answer Key}
\namedateperiod

\vspace{0.10in}
\begin{remindbox}
\small
\textbf{This is a practice test.} It mirrors the real Unit 1 test in length, format,
and the ideas it covers, but the numbers and contexts differ. Work every problem
with no notes, then check your answers against the key.
\end{remindbox}

% ======================================================================
\parthead{Part A --- Vocabulary (8 pts)}
Write the letter of the best description next to each term.

\vspace{4pt}
\noindent
\begin{minipage}[t]{0.44\textwidth}
\begin{enumerate}[leftmargin=2.2em, itemsep=7pt, label=\arabic*.]
  \item Scalar \ans{D}
  \item Magnitude \ans{F}
  \item Unit vector \ans{B}
  \item Linear combination \ans{G}
  \item Span \ans{H}
  \item Dot product \ans{E}
  \item Column space \ans{C}
  \item Rank \ans{A}
\end{enumerate}
\end{minipage}%
\hfill
\begin{minipage}[t]{0.53\textwidth}
\small
\begin{description}[leftmargin=1.4em, itemsep=3pt]
  \item[(A)] The number of independent columns of a matrix.
  \item[(B)] A vector whose length is exactly $1$.
  \item[(C)] All vectors you can reach as $A\xx$ --- the span of a matrix's columns.
  \item[(D)] A single number that scales (stretches or flips) a vector.
  \item[(E)] $\vv\cdot\ww=v_1w_1+v_2w_2$; it measures length and angle.
  \item[(F)] The length of a vector, $\sqrt{v_1^2+v_2^2}$.
  \item[(G)] A sum of scaled vectors, $c\vv+d\ww$.
  \item[(H)] The set of \emph{all} linear combinations of a set of vectors.
\end{description}
\end{minipage}

% ======================================================================
\parthead{Part B --- Multiple Choice (12 pts)}
Circle the best answer.

\begin{enumerate}[leftmargin=1.9em, itemsep=9pt]
  \item Two nonzero vectors are \textbf{perpendicular} exactly when
    \hfill (A) their sum is $\mathbf 0$ \quad \textbf{(B) $\vv\cdot\ww=0$} \ans{$\leftarrow$} \quad
    (C) $\vv\cdot\ww=1$ \quad (D) they have equal length
  \item The vectors $\begin{bsmallmatrix}1\\2\end{bsmallmatrix}$ and
    $\begin{bsmallmatrix}3\\6\end{bsmallmatrix}$ span
    \hfill (A) a single point \quad \textbf{(B) a line} \ans{$\leftarrow$} \quad (C) the whole plane \quad (D) all of space
  \item The product $A\xx$ is
    \hfill \textbf{(A) a combination of the \emph{columns} of $A$} \ans{$\leftarrow$} \quad (B) a combination of the rows of $A$ \\
    \phantom{x}\hfill (C) always $\mathbf 0$ \quad (D) the length of $\xx$
  \item The matrix $\begin{bsmallmatrix}1&2\\2&4\end{bsmallmatrix}$ has rank
    \hfill (A) $0$ \quad \textbf{(B) $1$} \ans{$\leftarrow$} \quad (C) $2$ \quad (D) $4$
  \item The unit vector $\dfrac{\vv}{\|\vv\|}$ has length
    \hfill (A) $0$ \quad \textbf{(B) $1$} \ans{$\leftarrow$} \quad (C) $\|\vv\|$ \quad (D) $\|\vv\|^2$
  \item A vector $\bb$ lies in the column space of $A$ exactly when
    \hfill (A) $\bb=\mathbf 0$ \quad \textbf{(B) $A\xx=\bb$ has a solution} \ans{$\leftarrow$} \\
    \phantom{x}\hfill (C) $\bb$ is a unit vector \quad (D) $\bb$ is perpendicular to every column
\end{enumerate}

% ======================================================================
\parthead{Part C --- Short Answer \& Computation (30 pts)}
Show your work.

\begin{enumerate}[leftmargin=1.9em, itemsep=12pt]
  \item Let $\vv=\begin{bsmallmatrix}3\\4\end{bsmallmatrix}$ and
        $\ww=\begin{bsmallmatrix}-1\\2\end{bsmallmatrix}$. Compute
        \textbf{(a)} $\vv+\ww$, \quad \textbf{(b)} $2\vv-\ww$, \quad \textbf{(c)} $\|\vv\|$.
        \begin{work}
          \vv+\ww &= \begin{bmatrix}3\\4\end{bmatrix}+\begin{bmatrix}-1\\2\end{bmatrix}
                   = \begin{bmatrix}2\\6\end{bmatrix} \\
          2\vv-\ww &= \begin{bmatrix}6\\8\end{bmatrix}-\begin{bmatrix}-1\\2\end{bmatrix}
                   = \begin{bmatrix}7\\6\end{bmatrix} \\
          \|\vv\| &= \sqrt{3^2+4^2} = \sqrt{25} = 5
        \end{work}
  \item Write the unit vector that points the same direction as
        $\vv=\begin{bsmallmatrix}3\\4\end{bsmallmatrix}$.
        \begin{work}
          \|\vv\| &= 5 \\
          \frac{\vv}{\|\vv\|} &= \frac{1}{5}\begin{bmatrix}3\\4\end{bmatrix}
                   = \begin{bmatrix}3/5\\4/5\end{bmatrix}
        \end{work}
  \item Let $\uu=\begin{bsmallmatrix}4\\2\end{bsmallmatrix}$ and
        $\vv=\begin{bsmallmatrix}-1\\2\end{bsmallmatrix}$. Compute $\uu\cdot\vv$, then state
        whether $\uu$ and $\vv$ are perpendicular and how you know.
        \begin{work}
          \uu\cdot\vv &= (4)(-1)+(2)(2) \\
                      &= -4+4 = 0 \\
                      &\Rightarrow\ \text{perpendicular --- the dot product is }0
        \end{work}
  \item Find weights $c$ and $d$ so that
        $c\begin{bsmallmatrix}2\\1\end{bsmallmatrix}+d\begin{bsmallmatrix}1\\3\end{bsmallmatrix}
         =\begin{bsmallmatrix}5\\5\end{bsmallmatrix}$.
        \begin{work}
          2c+d &= 5 \\
          c+3d &= 5 \\
          c &= 2, \quad d = 1 \\
          2\begin{bmatrix}2\\1\end{bmatrix}+1\begin{bmatrix}1\\3\end{bmatrix}
            &= \begin{bmatrix}5\\5\end{bmatrix}
        \end{work}
  \item Let $A=\begin{bsmallmatrix}1&2\\3&1\end{bsmallmatrix}$ and
        $\xx=\begin{bsmallmatrix}2\\1\end{bsmallmatrix}$. Compute $A\xx$ \emph{as a combination
        of the columns of $A$}.
        \begin{work}
          A\xx &= 2\begin{bmatrix}1\\3\end{bmatrix}+1\begin{bmatrix}2\\1\end{bmatrix} \\
               &= \begin{bmatrix}2\\6\end{bmatrix}+\begin{bmatrix}2\\1\end{bmatrix}
                = \begin{bmatrix}4\\7\end{bmatrix}
        \end{work}
  \item Factor $A=\begin{bsmallmatrix}1&2&3\\2&1&3\end{bsmallmatrix}$ as $A=CR$. Which columns are
        independent? Write $C$, build the recipe $R$, and give the rank.
        \begin{work}
          \text{col }3 &= \begin{bmatrix}1\\2\end{bmatrix}+\begin{bmatrix}2\\1\end{bmatrix}
                        = \begin{bmatrix}3\\3\end{bmatrix}
                        \quad \text{(columns }1,2\text{ independent)} \\
          C &= \begin{bmatrix}1&2\\2&1\end{bmatrix}, \quad
              R = \begin{bmatrix}1&0&1\\0&1&1\end{bmatrix} \\
          \text{rank} &= 2
        \end{work}
  \boxguard[14]
  \item Let $A=\begin{bsmallmatrix}1&2\\2&4\end{bsmallmatrix}$. For each target, say whether it is
        reachable as $A\xx$ (is it in the column space?) and justify in one line.
        \textbf{(a)} $\bb=\begin{bsmallmatrix}3\\6\end{bsmallmatrix}$ \quad
        \textbf{(b)} $\bb=\begin{bsmallmatrix}1\\1\end{bsmallmatrix}$
        \begin{work}
          C(A) &= \text{the line through } \begin{bmatrix}1\\2\end{bmatrix} \\
          \text{(a)}\quad \begin{bmatrix}3\\6\end{bmatrix}
            &= 3\begin{bmatrix}1\\2\end{bmatrix} \quad \text{reachable} \\
          \text{(b)}\quad \begin{bmatrix}1\\1\end{bmatrix}
            &\ne c\begin{bmatrix}1\\2\end{bmatrix} \quad \text{not reachable}
        \end{work}
\end{enumerate}

% ======================================================================
\parthead{Part D --- Extended Response (10 pts)}
Explain your reasoning in full sentences.

\begin{enumerate}[leftmargin=1.9em, itemsep=14pt]
  \item A matrix $A=\begin{bsmallmatrix}1&2\\1&3\end{bsmallmatrix}$ has columns
        $\begin{bsmallmatrix}1\\1\end{bsmallmatrix}$ and $\begin{bsmallmatrix}2\\3\end{bsmallmatrix}$.
        \textbf{(a)} Is the column space a line or the whole plane? How do you know?
        \textbf{(b)} Is $\bb=\begin{bsmallmatrix}4\\5\end{bsmallmatrix}$ reachable? If so, find the
        weights. \textbf{(c)} If the second column were changed to
        $\begin{bsmallmatrix}2\\2\end{bsmallmatrix}$, what would happen to the column space, and why?
        \ansline{(a) The whole plane --- the columns are not multiples, so they span $\mathbb R^2$.}
        \begin{work}
          c\begin{bmatrix}1\\1\end{bmatrix}+d\begin{bmatrix}2\\3\end{bmatrix}
            &= \begin{bmatrix}4\\5\end{bmatrix} \\
          c+2d &= 4 \\
          c+3d &= 5 \\
          d &= 1, \quad c = 2
        \end{work}
        \ansline{(c) Column $2$ becomes $2\times$ column $1$, so $C(A)$ shrinks to a line (rank $1$).}
  \item Two viewers' taste profiles are $\uu=\begin{bsmallmatrix}2\\1\end{bsmallmatrix}$ and
        $\vv=\begin{bsmallmatrix}-1\\2\end{bsmallmatrix}$. Compute $\uu\cdot\vv$ and explain what the
        result says about how similar the two viewers' tastes are, using what you know about the
        sign of a dot product.
        \begin{work}
          \uu\cdot\vv &= (2)(-1)+(1)(2) = 0
        \end{work}
        \ansline{Zero means perpendicular --- neither aligned nor opposed, so tastes are unrelated.}
        \ansline{A positive dot would mean similar tastes; a negative dot, opposite tastes.}
\end{enumerate}

\end{document}
